\chapter{Convexity and curvature pinching along the flow}\label{sec:mcf-pinching}
Throughout this section $F:M^n\times[0,T)\to\mathbb R^{n+1}$ is a smooth mean curvature flow of a closed manifold, with $n\geq2$, and with a smooth choice of unit normal. We retain $h(X,Y)=\langle D_X\nu,Y\rangle$, $A=h^\sharp$, $H=\operatorname{tr}A$, and $\partial_tF=-H\nu$. Weak convexity means $A\geq0$; strict convexity means $A>0$. Mean convexity $H\geq0$ is a weaker condition. The source is \cite{Hui84}, Sections 3--5, pp.~242--247. We prove the preservation and quotient estimates here; convergence of a normalized convex flow to a sphere requires further estimates.

\begin{lemma}[Preserving a nonnegative endomorphism with scalar reaction]\label{lem:evolving-metric-positive-tensor}
Let $g(t)$ be a smooth family of metrics on a closed manifold for $0\leq t\leq\tau$. Suppose $S(t)$ is a smooth endomorphism, self-adjoint with respect to $g(t)$, satisfying
\[
 (\partial_t-\Delta_{g(t)})S=aS
\]
for a smooth scalar function $a$. If $S(0)\geq0$, then $S(t)\geq0$ for every $t\leq\tau$. The Laplacian on endomorphisms uses the connection of $g(t)$. This elementary special case suffices for the tensor-preservation arguments of \cite{Hui84}, Section 4.
\end{lemma}
\begin{proof}
Choose $C>\sup a$ on the compact space-time interval and put $S_\delta=S+\delta e^{Ct}I$, $\delta>0$. It is positive definite initially and satisfies
\[
 (\partial_t-\Delta)S_\delta=aS_\delta+\delta e^{Ct}(C-a)I.
\]
If positivity first fails, compactness supplies a first time $t_0>0$, a point $x_0$, and a nonzero vector $v$ with $S_\delta v=0$. Extend $v$ smoothly near this point with spatial covariant derivative zero there, and hold its coordinate components fixed in time. Before the first failure the scalar $\psi=g(S_\delta v,v)$ is positive; at $(x_0,t_0)$ it is zero and has a spatial minimum. Thus $(\partial_t-\Delta)\psi\leq0$ there.

At that point the derivative of the metric and the second derivatives of $v$ contribute zero, because $S_\delta v=0$ and $S_\delta$ is self-adjoint; first spatial derivatives of $v$ are zero by construction. Consequently
\[
 (\partial_t-\Delta)\psi
 =g(((\partial_t-\Delta)S_\delta)v,v)
 =\delta e^{Ct_0}(C-a)|v|^2>0,
\]
a contradiction. Let $\delta\downarrow0$. The argument treats the evolving metric explicitly and does not assume a time-independent orthonormal frame.
\end{proof}

\begin{proposition}[Preservation of convexity and linear pinching]\label{prop:mcf-linear-pinching}
Along the flow above, $A\geq0$ is preserved. More generally, for any constants $\varepsilon,\beta\in\mathbb R$, each inequality
\[
 A\geq\varepsilon H I,\qquad A\leq\beta H I
\]
is preserved if it holds initially. If the initial hypersurface is strictly convex, then it remains strictly convex at every smooth time. Compare \cite{Hui84}, Corollary 4.2 and Theorem 4.3, p.~245.
\end{proposition}
\begin{proof}
Proposition \ref{prop:mcf-curvature-evolution} gives
$(\partial_t-\Delta)A=|A|^2A$ and $(\partial_t-\Delta)H=|A|^2H$. The identity endomorphism is parallel and has zero time derivative. Hence each of $A$, $A-\varepsilon HI$, and $\beta HI-A$ satisfies the scalar-reaction equation of Lemma \ref{lem:evolving-metric-positive-tensor}. Apply it on any compact smooth time interval.

If $A>0$ initially, then $H>0$ and compactness supplies $\varepsilon>0$ with $A\geq\varepsilon HI$ initially. Mean curvature stays positive by Corollary \ref{cor:mcf-mean-convexity}, applied on each connected component. The preserved lower bound now implies $A>0$. No assertion about evolution beyond the maximal smooth time is made.
\end{proof}

\begin{proposition}[Evolution of the normalized curvature ratio]\label{prop:mcf-curvature-ratio}
Wherever $H>0$, set $r=|A|^2/H^2$. Then
\[
 (\partial_t-\Delta)r
 =\frac2H\langle\nabla H,\nabla r\rangle
 -\frac{2}{H^4}|H\nabla A-A\otimes\nabla H|^2.
\]
This is the zero-parameter quotient identity in \cite{Hui84}, Lemma 5.2, pp.~246--247. The tensor in the last term has components $H\nabla_kh_{ij}-h_{ij}\nabla_kH$.
\end{proposition}
\begin{proof}
Write $Q=|A|^2$ and $L=\partial_t-\Delta$. The evolution identities yield $LQ=-2|\nabla A|^2+2Q^2$ and $LH=QH$. The scalar chain and product rules therefore give
\[
 L(QH^{-2})=-\frac{2|\nabla A|^2}{H^2}
 -\frac{6Q}{H^4}|\nabla H|^2
 +\frac4{H^3}\langle\nabla Q,\nabla H\rangle.
\]
The reaction terms cancel. Also
\[
 \frac2H\langle\nabla H,\nabla r\rangle
 =\frac2{H^3}\langle\nabla H,\nabla Q\rangle
 -\frac{4Q}{H^4}|\nabla H|^2.
\]
Subtract this expression. Expanding the square in the statement, and using $\nabla Q=2\langle A,\nabla A\rangle$, gives exactly the remaining three terms. The time-dependent metric has already been accounted for in the evolution of $Q$; treating the covariant components of $h$ as a fixed-metric matrix would not give that evolution correctly.
\end{proof}

\begin{corollary}[A nonincreasing curvature-ratio bound]\label{cor:mcf-ratio-bound}
If $H>0$ initially, then for every later smooth time
\[
 \frac1n\leq\frac{|A|^2}{H^2}\leq\max_{M_0}\frac{|A|^2}{H^2}.
\]
In particular an initial bound $|A|^2\leq CH^2$ is preserved, and the spatial maximum of $r$ is nonincreasing. This assertion does not by itself say that the maximum converges to $1/n$ at a singular time.
\end{corollary}
\begin{proof}
Positivity of $H$ persists. On each compact smooth time interval, $H$ has a positive minimum, so the drift $2\nabla H/H$ in Proposition \ref{prop:mcf-curvature-ratio} is smooth and bounded. The scalar parabolic maximum principle, applied from any starting time in the interval, makes the maximum of $r$ nonincreasing. The lower bound is the pointwise Cauchy--Schwarz inequality $(\sum_i\kappa_i)^2\leq n\sum_i\kappa_i^2$. Equality in this algebraic inequality means all principal curvatures agree. Compare \cite{Hui84}, Section 5; improved decay of the trace-free ratio requires more than this maximum-principle bound.
\end{proof}

\begin{lemma}[Parallel normalized curvature under strict convexity]\label{lem:parallel-normalized-shape}
Let a connected Euclidean hypersurface of dimension $n\geq2$ be strictly convex and satisfy $\nabla(A/H)=0$. Then it is contained in a round sphere. If it is compact without boundary, its image is the entire sphere.
\end{lemma}
\begin{proof}
At a point use a principal orthonormal frame. The hypothesis says $\nabla_i h_{jk}=h_{jk}\nabla_iH/H$. Codazzi gives
$h_{jk}\nabla_iH=h_{ik}\nabla_jH$. For each $i$ choose $j\ne i$ and put $k=j$; strict positivity of $\kappa_j$ forces $\nabla_iH=0$. Thus $H$ is constant and $\nabla h=0$.

Since $h$ is parallel, its derivative commutator is zero. Lemma \ref{lem:two-tensor-commutation}, applied to $e_i,e_j,e_i,e_j$, yields $K_{ij}(\kappa_i-\kappa_j)=0$. The Euclidean Gauss equation gives $K_{ij}=\kappa_i\kappa_j>0$, so all eigenvalues agree: $A=\lambda I$ with $\lambda=H/n>0$ constant. Differentiating $F-\nu/\lambda$ now gives zero. Connectedness makes this vector a fixed center, and $|F-\text{center}|=1/\lambda$. The image of a compact immersed manifold of the same dimension as this sphere is open by the inverse function theorem and closed by compactness, hence is the full connected sphere. No global equality-classification theorem is used.
\end{proof}

\begin{theorem}[Rigidity when the maximum curvature ratio stops decreasing]\label{thm:mcf-ratio-rigidity}
Suppose $M$ is connected and the initial hypersurface is strictly convex. If for some smooth time $t_0>0$,
\[
 \max_{M_{t_0}}r=\max_{M_0}r,
\]
then every slice for $0\leq t\leq t_0$ has image a round sphere, and $r=1/n$ there. This is a consequence of the quotient identity in \cite{Hui84}, Lemma 5.2, together with the strong scalar maximum principle; it is distinct from the paper's general convergence theorem.
\end{theorem}
\begin{proof}
Let $C=\max_{M_0}r$. Proposition \ref{prop:mcf-linear-pinching} preserves strict convexity, and Corollary \ref{cor:mcf-ratio-bound} gives $r\leq C$. At time $t_0$, the nonnegative function $C-r$ vanishes somewhere. The strong parabolic maximum principle for the drift operator in Proposition \ref{prop:mcf-curvature-ratio} makes it zero on $M\times[0,t_0]$. To treat the top time as an interior time one can work on a slightly longer compact smooth interval, which exists because $t_0<T$; the conclusion propagates backward to the initial time by continuity.

The quotient identity then forces $H\nabla A-A\otimes\nabla H=0$ on each slice at positive time, which is $\nabla(A/H)=0$. Lemma \ref{lem:parallel-normalized-shape} makes each such slice a round sphere, where $r=1/n$. Continuity gives the same identity and umbilicity at time zero; strict positivity and the lemma give the initial sphere as well. The scalar strong maximum principle is the analytic background input in this argument.
\end{proof}

\begin{example}[Cylinders have constant ratio without being umbilic]\label{ex:cylinder-curvature-ratio}
For $1\leq k<n$, the shrinking round cylinder $\mathbb S^k(R(t))\times\mathbb R^{n-k}$ with $R(t)^2=R(0)^2-2kt$ has
\[
 H=\frac{k}{R(t)},\qquad |A|^2=\frac{k}{R(t)^2},\qquad r=\frac1k>\frac1n.
\]
It has positive mean curvature and a constant curvature ratio, but is neither strictly convex nor compact. Thus it does not meet the hypotheses of Theorem \ref{thm:mcf-ratio-rigidity}.
\end{example}
\begin{proof}
The $k$ spherical principal curvatures are $1/R(t)$; the remaining $n-k$ principal curvatures are zero. Summing them and their squares gives the formulas. The radius equation $R'=-k/R$ is precisely the inward normal mean-curvature velocity. On each product slice the second fundamental form and $H$ are spatially parallel, so the negative square in Proposition \ref{prop:mcf-curvature-ratio} vanishes. The ratio remains $1/k$ because both numerator and denominator scale as $R^{-2}$. The zero principal curvatures explain exactly where the strict-convexity argument in Lemma \ref{lem:parallel-normalized-shape} fails.
\end{proof}
